\textbf{Phần 1. Trắc nghiệm lựa chọn}
\setcounterpageref{ex}{0}

\begin{ex} %[4D8V44]
Kết quả của phép chia $288 : 24$ là
\choice
{$11$}
{$13$}
{\True $12$}
{$14$}
\loigiai{
$288 : 24 = 12$.\\
\textbf{Đáp án của câu hỏi là C}
}
\end{ex}

\begin{ex} %[4D8V44]
Số dư trong phép chia $150 : 20$ là
\choice
{$0$}
{\True $10$}
{$5$}
{$15$}
\loigiai{
$150 : 20 = 7$ dư 10.\\
\textbf{Đáp án của câu hỏi là B}
}
\end{ex}

\begin{ex} %[4D8V44]
Kết quả của phép chia $1440 : 36$ là
\choice
{$30$}
{$50$}
{$45$}
{\True $40$}
\loigiai{
$1440 : 36 = 40$.\\
\textbf{Đáp án của câu hỏi là D}
}
\end{ex}

\begin{ex} %[4D8V44]
Tìm $x$ biết: $x \times 25 = 1250$.
\choice
{\True $50$}
{$40$}
{$60$}
{$55$}
\loigiai{
$x = 1250 : 25 = 50$.\\
\textbf{Đáp án của câu hỏi là A}
}
\end{ex}

\begin{ex} %[4D8V44]
Một hình chữ nhật có diện tích $360\text{ m}^2$, chiều rộng $15\text{ m}$. Chiều dài của hình chữ nhật là
\choice
{$22\text{ m}$}
{$26\text{ m}$}
{\True $24\text{ m}$}
{$28\text{ m}$}
\loigiai{
$360 : 15 = 24\text{ m}$.\\
\textbf{Đáp án của câu hỏi là C}
}
\end{ex}

\begin{ex} %[4D8V44]
Kết quả của phép chia $7200 : 45$ là
\choice
{$150$}
{\True $160$}
{$170$}
{$180$}
\loigiai{
$7200 : 45 = 160$.\\
\textbf{Đáp án của câu hỏi là B}
}
\end{ex}

\begin{ex} %[4D8V44]
Người ta xếp $1800$ quyển vở vào các thùng, mỗi thùng chứa được $40$ quyển. Số thùng cần dùng là
\choice
{$40$ thùng}
{$50$ thùng}
{$55$ thùng}
{\True $45$ thùng}
\loigiai{
$1800 : 40 = 45$ thùng.\\
\textbf{Đáp án của câu hỏi là D}
}
\end{ex}

\begin{ex} %[4D8V44]
Một đoàn khách du lịch gồm $265$ người. Người ta thuê các xe ô tô loại $29$ chỗ ngồi để đi tham quan. Cần thuê ít nhất bao nhiêu xe như thế?
\choice
{\True $10$ xe}
{$9$ xe}
{$11$ xe}
{$12$ xe}
\loigiai{
$265 : 29 = 9$ dư 4 người. Cần ít nhất $9 + 1 = 10$ xe.\\
\textbf{Đáp án của câu hỏi là A}
}
\end{ex}

\begin{ex} %[4D8V44]
Tính giá trị biểu thức: $4500 : 15 + 3600 : 18$.
\choice
{$400$}
{$600$}
{\True $500$}
{$550$}
\loigiai{
$300 + 200 = 500$.\\
\textbf{Đáp án của câu hỏi là C}
}
\end{ex}

\begin{ex} %[4D8V44]
Một xí nghiệp có $3600\text{ kg}$ đường đóng vào các bao loại $50\text{ kg}$. Người ta đã bán đi $40$ bao đường đó. Số bao đường còn lại là
\choice
{$30$ bao}
{\True $32$ bao}
{$35$ bao}
{$40$ bao}
\loigiai{
Tổng số bao: $3600 : 50 = 72$ bao.\\
Còn lại: $72 - 40 = 32$ bao.\\
\textbf{Đáp án của câu hỏi là B}
}
\end{ex}

\begin{ex} %[4D8V44]
Một mảnh đất hình chữ nhật có chu vi $140\text{ m}$, chiều rộng $25\text{ m}$. Diện tích mảnh đất là
\choice
{$1000\text{ m}^2$}
{$1200\text{ m}^2$}
{$1250\text{ m}^2$}
{\True $1125\text{ m}^2$}
\loigiai{
Nửa chu vi: $70\text{ m}$. Chiều dài: $70 - 25 = 45\text{ m}$.\\
Diện tích: $45 \times 25 = 1125\text{ m}^2$.\\
\textbf{Đáp án của câu hỏi là D}
}
\end{ex}

\begin{ex} %[4D8V44]
Tìm số tự nhiên $x$ lớn nhất có hai chữ số sao cho $2500 : x > 50$.
\choice
{\True $49$}
{$48$}
{$50$}
{$51$}
\loigiai{
$2500 : 50 = 50 \Rightarrow x < 50$. Số tự nhiên $x$ lớn nhất có hai chữ số là 49.\\
\textbf{Đáp án của câu hỏi là A}
}
\end{ex}

\textbf{Phần 2. Trắc nghiệm đúng sai}
\setcounterpageref{ex}{0}

\begin{ex} %[4D8V44]
Xét tính đúng sai của các phép chia cho số có hai chữ số:
\choiceTF
{\True $480 : 12 = 40$.}
{$750 : 25 = 25$.}
{\True $750 : 25 = 30$.}
{\True Trong phép chia có dư, số dư luôn nhỏ hơn số chia.}
\loigiai{
\begin{itemize}
    \item $480 : 12 = 40$. Đúng.
    \item $750 : 25 = 30 \neq 25$. Ý 2 Sai, Ý 3 Đúng.
    \item Số dư $<$ số chia: Đúng.
\end{itemize}
\begin{itemize}
    \item a) \textbf{Đúng}.
    \item b) \textbf{Sai}.
    \item c) \textbf{Đúng}.
    \item d) \textbf{Đúng}.
\end{itemize}
}
\end{ex}

\begin{ex} %[4D8V44]
Một cửa hàng nhập về $1200\text{ kg}$ gạo tẻ và $800\text{ kg}$ gạo nếp đóng vào các bao loại $40\text{ kg}$. Xét tính đúng sai:
\choiceTF
{\True Số bao gạo tẻ là $30$ bao.}
{\True Số bao gạo nếp là $20$ bao.}
{Tổng số bao gạo là $60$ bao.}
{\True Tổng số bao gạo là $50$ bao.}
\loigiai{
\begin{itemize}
    \item Gạo tẻ: $1200 : 40 = 30$ bao. Đúng.
    \item Gạo nếp: $800 : 40 = 20$ bao. Đúng.
    \item Tổng: $30 + 20 = 50$ bao $\neq 60$ bao. Ý 3 Sai, Ý 4 Đúng.
\end{itemize}
\begin{itemize}
    \item a) \textbf{Đúng}.
    \item b) \textbf{Đúng}.
    \item c) \textbf{Sai}.
    \item d) \textbf{Đúng}.
\end{itemize}
}
\end{ex}

\textbf{Phần 3. Trả lời ngắn}
\setcounterpageref{ex}{0}

\begin{ex} %[4D8V44]
Tính thương: $720 : 18$.
\shortans{40}
\loigiai{
\textbf{Đáp số:} $40$\\
$40$.
}
\end{ex}

\begin{ex} %[4D8V44]
Tính thương: $1250 : 25$.
\shortans{50}
\loigiai{
\textbf{Đáp số:} $50$\\
$50$.
}
\end{ex}

\begin{ex} %[4D8V44]
Tìm $x$ biết: $x \times 35 = 1400$.
\shortans{40}
\loigiai{
\textbf{Đáp số:} $40$\\
$x = 1400 : 35 = 40$.
}
\end{ex}

\begin{ex} %[4D8V44]
Một mảnh vườn hình chữ nhật có diện tích $480\text{ m}^2$, chiều dài $24\text{ m}$. Tính chu vi mảnh vườn theo mét.
\shortans{88}
\loigiai{
\textbf{Đáp số:} $88$\\
Rộng: $480 : 24 = 20\text{ m}$. Chu vi: $(24 + 20) \times 2 = 88\text{ m}$.
}
\end{ex}

\begin{ex} %[4D8V44]
Có $960$ quyển sách xếp đều vào 24 ngăn tủ. Hỏi mỗi ngăn tủ có bao nhiêu quyển sách?
\shortans{40}
\loigiai{
\textbf{Đáp số:} $40$\\
$960 : 24 = 40$.
}
\end{ex}

\begin{ex} %[4D8V44]
Tìm số bị chia trong phép chia cho 32 có thương là 45 và số dư là số dư lớn nhất có thể.
\shortans{1471}
\loigiai{
\textbf{Đáp số:} $1471$\\
Số dư lớn nhất: 31. Số bị chia: $45 \times 32 + 31 = 1440 + 31 = 1471$.
}
\end{ex}
